% !TEX root =  ../master.tex

% Abkürzungen
\newacronym{aos}{AoS}{Array of Structures}
\newacronym{api}{API}{Application Programming Interface}
\newacronym{ast_acr}{AST}{Abstract Syntax Tree}
\newacronym{cpu}{CPU}{Central Processing Unit}
\newacronym{cst_acr}{CST}{Concrete Syntax Tree}
\newacronym{dfa}{DFA}{Deterministic Finite Automaton}
\newacronym{ebnf}{EBNF}{Erweiterte Backus-Naur-Form}
\newacronym{gc}{GC}{Garbage Collection}
\newacronym{ide}{IDE}{Integrated Development Environment}
\newacronym{ieee}{IEEE}{Institute of Electrical and Electronics Engineers}
\newacronym{io}{I/O}{Input/Output}
\newacronym{json}{JSON}{JavaScript Object Notation}
\newacronym{llm}{LLM}{Large Language Model}
\newacronym{nfa}{NFA}{Nondeterministic Finite Automaton}
\newacronym{ply}{PLY}{Python Lex-Yacc}
\newacronym{ram}{RAM}{Random Access Memory}
\newacronym{soa_acr}{SoA}{Structure of Arrays}
\newacronym{utf}{UTF}{Unicode Transformation Format}
\newacronym{v8}{V8}{JavaScript-Engine (Google)}


% Glossareinträge
\newglossaryentry{amortisiert}{
    name={Amortisierte Laufzeit},
    description={Durchschnittlicher Aufwand einer Operation über eine sehr lange Sequenz von Aufrufen, wodurch seltene kostenintensive Operationen rechnerisch auf viele effiziente Operationen umgelegt werden}
}

\newglossaryentry{backshift}{
    name={Backshift Deletion},
    description={Verfahren zum Löschen in Hashtabellen mit Open Addressing, bei dem nachfolgende Elemente aufrücken, um Lücken zu schließen und die Sondierungskette ohne Tombstones zu reparieren}
}

\newglossaryentry{cst}{
    name={Concrete Syntax Tree (CST)},
    description={Vollständiger Ableitungsbaum einer Grammatik, der alle syntaktischen Details einschließlich Hilfszeichen wie Klammern repräsentiert}
}

\newglossaryentry{ast}{
    name={Abstract Syntax Tree (AST)},
    description={Verdichtete Baumdarstellung, die nur die für die Programmlogik wesentlichen Informationen wie Operatoren und Operanden enthält}
}

\newglossaryentry{garbagecollection}{
    name={Garbage Collection (GC)},
    description={Automatischer Prozess der Laufzeitumgebung zur Ermittlung und Freigabe von nicht mehr referenziertem Arbeitsspeicher}
}

\newglossaryentry{inkrementell}{
    name={Inkrementelles Parsing},
    description={Analyseverfahren, bei dem nach einer Änderung am Quelltext nur der betroffene Teil des Syntaxbaums neu geparst wird, anstatt das gesamte Dokument}
}

\newglossaryentry{linearprobing}{
    name={Linear Probing},
    description={Verfahren zur Kollisionsauflösung in Hashtabellen, bei dem bei einem besetzten Slot linear der nächste freie Speicherplatz gesucht wird}
}

\newglossaryentry{listenfaltung}{
    name={Listen-Faltung},
    description={Umwandlung einer flachen Folge von Elementen in eine rekursiv verschachtelte Listenstruktur (Cons-Zellen)}
}

\newglossaryentry{structuraltyping}{
    name={Strukturelles Typing},
    description={Konzept in TypeScript (statisches Duck-Typing), bei dem die Typkompatibilität rein anhand der Form eines Objekts geprüft wird}
}

\newglossaryentry{soa}{
    name={Structure of Arrays (SoA)},
    description={Architekturmuster, bei dem Objekteigenschaften in getrennten, parallelen Arrays gespeichert werden, um die Cache-Lokalität der CPU zu erhöhen}
}

\newglossaryentry{tombstone}{
    name={Tombstone},
    description={Marker-Objekt in einer Hashtabelle, das einen gelöschten Eintrag ersetzt, um ein Abreißen der Suchkette beim Linear Probing zu verhindern}
}

\newglossaryentry{typeassertion}{
    name={Type Assertion},
    description={Compiler-Direktive (as), die TypeScript zwingt, einen Wert als bestimmten Typ zu behandeln, ohne den Wert zur Laufzeit tatsächlich zu konvertieren}
}

\newglossaryentry{typenarrowing}{
    name={Type Narrowing},
    description={Prozess der schrittweisen Typ-Einengung durch den Compiler auf Basis von Kontrollflussanalysen wie typeof-Prüfungen oder if-Ver\-zwei\-gungen}
}